\chapter{Analyse und Konzeption}

\section{Verwandte Arbeiten}

Anbei werden Arbeiten aufgelistet, die Ähnlichkeiten zu dieser Arbeit haben und wichtige Vorarbeit geleistet haben:

\subsection{Flowable}

Flowable ist eine "Workflow Management Platform und Java Geschäftsprozessengine", welche Notationen wie BPMN, CMMN und DMN 
unterstützt und ausführt. \cite{doan_using_2022} Auch ist es möglich Flowable als Engine in eigenen Java-Projekten zu nutzen. \cite{flowable_ag_getting_nodate}
Auch wird damit geworben, dass Flowable eine IoT-Anbindung anbietet. Auf diese wird in dessen Produktdokumentation nicht eingegangen. Somit ist 
die Anbindung von IoT Systemen bei der Flowable Prozessengine nicht nachweisbar. \cite{flowable_ag_digitize_nodate}

\subsection{Camunda}

Bei Camunda handelt es sich ebenfalls um eine Workflow Engine, welche verwendet wird, um Geschäftsprozesse, welche mit der BPMN 2.0 Notation
modelliert worden sind, auszuführen. \cite{david_exploring_2023} Die Software wurde mit Java entwickelt und die Prozessengine läuft 
dementsprechend auf der Java Virtual Machine. Darüber hinaus hat sie auch eine Integration zu den Frameworks Spring und Java EE, sodass Camunda auch
als Bibliothek von Entwicklern genutzt werden kann. \cite{fernandez_camunda_2013}

\subsection{EMD - Erweiterbare und adaptive Elektro-Mobilitätsdienste}

In dieser Arbeit geht es um die Entwicklung von Diensten und Werkzeugen, um Innovationen im E-Mobilitätsbereich 
zu erforschen. Ein Produkt aus dieser Arbeit ist ein weiterentwickeltes BPMN Tool, das benötigte Dienste integrieren und 
Prozesse durch einen BPMN-Interpreter ausführen und orchestrieren kann. Dabei wurden Geräte wie eBike-Ladestationen,
Parkplatzschranke und gerätenahe Sensorik verwendet.\cite{fahndrich_emd_2016}

\subsection{Agent-based Business Process Orchestration for IoT}

In dieser Arbeit wird ein Geschäftsprozess in BPMN modelliert und mit Agenten und IoT verknüpft. Dabei kommen auch menschliche
Akteure ins Spiel. Standards die dabei verwendet werden sind O-MI und O-DF. Dabei wurde aber nicht auf echter Hardware, sondern mit 
Mock-Ups gearbeitet. Technologien die für die Orchestrierungsebene verwendet worden sind, sind Java und Node.js. 
In der Diskussion wurde die Embedded-Schicht als Thema für zukünftige Forschungsarbeiten aufgegriffen. \cite{kampik_agent-based_2019}

\subsection{Bridging BPMN and Edge-IoT: A Reactor-based architecture for real-time, AI-driven workflow orchestration}

Hier wurde eine Reactor-basierte Architektur entwickelt, um Edge-Geräte durch die Zeebe-Execution-Engine des Herstellers Camunda auf 
BPMN Prozesse reagieren zu lassen. Dies wurde noch nicht auf richtiger Hardware ausgeführt, sondern mit Simulationen. Dabei dient die Zeebe-Engine
als Orchestrator. \cite{thahir_bridging_2026}

\subsection{Modelling and executing IoT-enhanced business processes through BPMN and microservices}

In der Arbeit \textit{Modelling and executing IoT-enhanced business processes through BPMN and microservices} behandelt Valderas die Ausführung von 
IoT-gekoppelten Prozessen. Dabei werden die Prozesse in BPMN notiert und mit der Camunda Execution Engine ausgeführt, welche einen Raspberry Pi ansteuern. 
\cite{valderas_modelling_2022}

\subsection{Evaluation der verwandten Arbeiten}

All diese Arbeiten haben die Modellierungssprache BPMN oder CMMN verwendet, jedoch wurden EPC's in diesem Kontext noch nicht beleuchtet. Und das ist der Beitrag dieser
Arbeit. Dieser beinhaltet einen Compiler, welcher EPC-Modelle in ein ausführbares Format umwandelt und auf dieser Basis IoT-Geräte über das CoAP Protokoll ansteuert.
Darüber hinaus wird auch der Rahmen für Firmware-Code damit generiert, sodass die Entwicklung von prozessorientierter Software erleichtert werden kann.

\section{Systemanforderungen}

Die Systemanforderungen kann man anhand von User Stories definieren. Diese werden nach dem Verfahren nach Cohn \cite{cohn_user_2004} aufgezählt:

\begin{enumerate}[label=\Alph*)]
    \item Ich als \textbf{Geschäftsleiter} möchte meine modellierten \textbf{Prozesse auf Produktionshardware ausführen und orchestrieren können}, 
    sodass \textbf{Abläufe gemäß der Prozessdefinition laufen und das Unternehmen von prozessorientierter Organisation profitieren kann}.

    \item Ich als \textbf{Controller} möchte \textbf{laufende Prozesse überwachen können}, sodass \textbf{die Effizienz von Geschäftsprozessen gemessen
    und optimiert werden kann}. 

    \item Ich als \textbf{Softwareentwickler} möchte eine Programmierschnittstelle zur Prozessschicht, sodass \textbf{Produktionshardware leicht an das Prozessgeschehen
    angebunden werden kann}.

    \item Ich als \textbf{Prozessanalyst} möchte \textbf{einen Verlauf vergangener Prozesse} einsehen, sodass \textbf{Schwachstellen und Engpässe bei Prozessen ermittelt werden können}.
\end{enumerate}